% Options for packages loaded elsewhere
\PassOptionsToPackage{unicode}{hyperref}
\PassOptionsToPackage{hyphens}{url}
\documentclass[
  12pt,
]{ctexart}
\usepackage{xcolor}
\usepackage{amsmath,amssymb}
\setcounter{secnumdepth}{-\maxdimen} % remove section numbering
\usepackage{iftex}
\ifPDFTeX
  \usepackage[T1]{fontenc}
  \usepackage[utf8]{inputenc}
  \usepackage{textcomp} % provide euro and other symbols
\else % if luatex or xetex
  \usepackage{unicode-math} % this also loads fontspec
  \defaultfontfeatures{Scale=MatchLowercase}
  \defaultfontfeatures[\rmfamily]{Ligatures=TeX,Scale=1}
\fi
\usepackage{lmodern}
\ifPDFTeX\else
  % xetex/luatex font selection
\fi
% Use upquote if available, for straight quotes in verbatim environments
\IfFileExists{upquote.sty}{\usepackage{upquote}}{}
\IfFileExists{microtype.sty}{% use microtype if available
  \usepackage[]{microtype}
  \UseMicrotypeSet[protrusion]{basicmath} % disable protrusion for tt fonts
}{}
\makeatletter
\@ifundefined{KOMAClassName}{% if non-KOMA class
  \IfFileExists{parskip.sty}{%
    \usepackage{parskip}
  }{% else
    \setlength{\parindent}{0pt}
    \setlength{\parskip}{6pt plus 2pt minus 1pt}}
}{% if KOMA class
  \KOMAoptions{parskip=half}}
\makeatother
\usepackage{longtable,booktabs,array}
\usepackage{calc} % for calculating minipage widths
% Correct order of tables after \paragraph or \subparagraph
\usepackage{etoolbox}
\makeatletter
\patchcmd\longtable{\par}{\if@noskipsec\mbox{}\fi\par}{}{}
\makeatother
% Allow footnotes in longtable head/foot
\IfFileExists{footnotehyper.sty}{\usepackage{footnotehyper}}{\usepackage{footnote}}
\makesavenoteenv{longtable}
\setlength{\emergencystretch}{3em} % prevent overfull lines
\providecommand{\tightlist}{%
  \setlength{\itemsep}{0pt}\setlength{\parskip}{0pt}}
\usepackage{bookmark}
\IfFileExists{xurl.sty}{\usepackage{xurl}}{} % add URL line breaks if available
\urlstyle{same}
\hypersetup{
  hidelinks,
  pdfcreator={LaTeX via pandoc}}

\author{SCX}
\date{2026}

\begin{document}

\section{SCX 论文FrameworkReviewReport ---
五Review员Synthesis}\label{scx-ux8bbaux6587ux6846ux67b6ux5ba1ux67e5ux62a5ux544a-ux4e94ux5ba1ux67e5ux5458ux7efcux5408}

\begin{quote}
Review日期：2026-06-26 Review方式：5 个Independent AI Review员并行Review + Synthesis裁决
Review员Role：叙事弧线 / 理论构建 / Evidence链 / CompetitionPositioning / 跨DomainExtension
\end{quote}

\begin{center}\rule{0.5\linewidth}{0.5pt}\end{center}

\subsection{一、五Review员ConsensusMatrix}\label{ux4e00ux4e94ux5ba1ux67e5ux5458ux5171ux8bc6ux77e9ux9635}

\begin{longtable}[]{@{}lcccccc@{}}
\toprule\noalign{}
JudgmentDimension & 叙事 & 理论 & Evidence & Competition & 跨Domain & Consensus强度 \\
\midrule\noalign{}
\endhead
\bottomrule\noalign{}
\endlastfoot
EGP P1 先发是正确的 & ✅ & ✅ & ✅ & ✅ & ✅ & \textbf{全票} \\
SCX-Theory Independent投稿高风险 & ✅ & ✅ & ✅ & ✅ & ✅ & \textbf{全票} \\
V(s) 乘法分解是理论最薄弱环节 & --- & ✅ & --- & ✅ & --- & 2/2 评审 \\
``预估 29-48\%'' 需要重训Verification & ✅ & --- & ✅ & --- & --- & 2/2 评审 \\
应合并 Theory + MLIP 为一篇 & ✅ & --- & --- & --- & --- & 部分Consensus \\
MedMNIST/CIFAR Experiment当前有害 & --- & --- & ✅ & --- & ✅ & 2/2 评审 \\
``六Dimension覆盖''应改为单核心叙事 & ✅ & --- & --- & ✅ & ✅ & 3/3 评审 \\
Arrow 类比应删除 & --- & ✅ & --- & --- & --- & 单审 \\
与 Data Shapley 关系应为互补 & --- & --- & --- & ✅ & --- & 单审 \\
Domain从 7 个收窄到 1-2 个 & --- & --- & --- & --- & ✅ & 单审 \\
\end{longtable}

\begin{center}\rule{0.5\linewidth}{0.5pt}\end{center}

\subsection{二、核心恐惧Diagnosis：``别人觉得我是骗子''}\label{ux4e8cux6838ux5fc3ux6050ux60e7ux8bcaux65adux522bux4ebaux89c9ux5f97ux6211ux662fux9a97ux5b50}

\subsubsection{2.1
恐惧是否合理？}\label{ux6050ux60e7ux662fux5426ux5408ux7406}

\textbf{合理。} 不是因为 SCX
是假的，而是因为叙事顺序决定了审稿人的第一印象。

\subsubsection{2.2
三条攻击路径}\label{ux4e09ux6761ux653bux51fbux8defux5f84}

\paragraph{攻击路径
1：平凡性攻击（``这不就是条件概率吗''）}\label{ux653bux51fbux8defux5f84-1ux5e73ux51e1ux6027ux653bux51fbux8fd9ux4e0dux5c31ux662fux6761ux4ef6ux6982ux7387ux5417}

\begin{itemize}
\tightlist
\item
  SCX 的三个Proposition中，Prop 1 在数学上几乎 trivial
\item
  审稿人：``Conditional risk is a concept from Kolmogorov (1933). What's
  new here?''
\item
  \textbf{防御}：Proof\textbf{在真实材料体系中，忽略State条件性的代价有多大}------EGP
  P1 的 9.7 倍 per-batch RMSE 差异
\end{itemize}

\paragraph{攻击路径
2：拼凑攻击（``聚类+Model选择+主动学习拼在一起起了个新名字''）}\label{ux653bux51fbux8defux5f84-2ux62fcux51d1ux653bux51fbux805aux7c7bux6a21ux578bux9009ux62e9ux4e3bux52a8ux5b66ux4e60ux62fcux5728ux4e00ux8d77ux8d77ux4e86ux4e2aux65b0ux540dux5b57}

\begin{itemize}
\tightlist
\item
  审稿人：``State discovery = clustering. Expert reliability =
  per-cluster accuracy. Data value = heuristic product.''
\item
  \textbf{防御}：两层Descriptor（ErrorDrivenEncoder）是Core------F1 从 0.253
  → 0.585 的 2.3 倍提升Proof这不是 trivial clustering
\end{itemize}

\paragraph{攻击路径
3：空中楼阁攻击（``没有真Experiment证''）}\label{ux653bux51fbux8defux5f84-3ux7a7aux4e2dux697cux9601ux653bux51fbux6ca1ux6709ux771fux5b9eux9a8cux8bc1}

\begin{itemize}
\tightlist
\item
  审稿人：``You estimate 29-48\% improvement but never actually
  retrained.''
\item
  \textbf{防御}：一次重训（4-8
  小时计算）就能堵上------这是当前最致命的缺口
\end{itemize}

\begin{center}\rule{0.5\linewidth}{0.5pt}\end{center}

\subsection{三、最优论文Framework：``三段跳''}\label{ux4e09ux6700ux4f18ux8bbaux6587ux6846ux67b6ux4e09ux6bb5ux8df3}

\subsubsection{3.1 总体策略}\label{ux603bux4f53ux7b56ux7565}

\begin{verbatim}
EGP P1 (Independent, concrete)  →  SCX-MLIP (理论+应用合并, 核心战场)  →  SCX-Sim 或 SCX-Health (选一个深做)
\end{verbatim}

核心逻辑：\textbf{先 concrete 打根据地 → 再 abstract
建理论（但不Independent投稿）→ 再跨DomainExtension}

\begin{center}\rule{0.5\linewidth}{0.5pt}\end{center}

\subsubsection{3.2 第一跳：EGP Paper
1}\label{ux7b2cux4e00ux8df3egp-paper-1}

\begin{longtable}[]{@{}
  >{\raggedright\arraybackslash}p{(\linewidth - 2\tabcolsep) * \real{0.5000}}
  >{\raggedright\arraybackslash}p{(\linewidth - 2\tabcolsep) * \real{0.5000}}@{}}
\toprule\noalign{}
\begin{minipage}[b]{\linewidth}\raggedright
Dimension
\end{minipage} & \begin{minipage}[b]{\linewidth}\raggedright
说明
\end{minipage} \\
\midrule\noalign{}
\endhead
\bottomrule\noalign{}
\endlastfoot
\textbf{核心Contribution} & gauge-fixed shared+correction ACE Parameter化 \\
\textbf{Goal期刊} & npj Computational Materials / Physical Review
Materials / JCTC \\
\textbf{受众} & 材料/计算化学社区 \\
\textbf{篇幅} & 4-6 页正文 + SI \\
\textbf{Data} & AlN v3: 425-534 frames DFT 已Training; GaN/AlN v4
超算运行中 \\
\end{longtable}

\paragraph{章节结构}\label{ux7ae0ux8282ux7ed3ux6784}

\begin{verbatim}
1. Introduction
   - 多元素 ACE 势Function的 gauge ambiguity Problem
   - 已有Method（single ACE Training后拼接）的局限
   - 本文Contribution：shared+correction Parameter化 + coefficient-level gauge fixing

2. Methods
   2.1 Shared+Correction ACE Architecture
       E = Σ_i [c₀ᵀB_i + c_{Z_i}ᵀB_i + b_{Z_i}]
   2.2 Gauge Freedom and Post-Hoc Projection
       c_Z → c_Z - g, c₀ → c₀ + g (g = Σ π_Z c_Z)
   2.3 Why Soft Constraint Fails (gradient competition analysis)
   2.4 Training Protocol

3. Results: AlN v3
   3.1 EOS (V₀, B₀ vs DFT)
   3.2 Elastic Constants (C₁₁-C₆₆)
   3.3 Phonon Forces
   3.4 Gauge Violation: 8.77 → 4.6×10⁻¹⁶

4. Discussion
   - Gauge fixing 对势Function合并的意义（AlN+GaN→AlGaN）
   - 对纯 AlN 而言 gauge fixing 是形式完备性
   - Future work: state-conditioned expert reliability（一句话）
\end{verbatim}

\paragraph{红线（不可逾越）}\label{ux7ea2ux7ebfux4e0dux53efux903eux8d8a}

\begin{itemize}
\tightlist
\item
  ❌ 不塞 SCX 理论（最多 Discussion 中一句 future work）
\item
  ❌ 不包含 SCX 噪声检测 Figs（Fig 4-8 留给 SCX-MLIP）
\item
  ❌ 不声称 Model B 是''新的运行时 MoE 势FunctionArchitecture''
\item
  ❌ 不声称对 AlGaN 可转移（除非有ExperimentData）
\end{itemize}

\begin{center}\rule{0.5\linewidth}{0.5pt}\end{center}

\subsubsection{3.3
第二跳：SCX-MLIP（理论+应用合并，核心战场）}\label{ux7b2cux4e8cux8df3scx-mlipux7406ux8bbaux5e94ux7528ux5408ux5e76ux6838ux5fc3ux6218ux573a}

\begin{longtable}[]{@{}ll@{}}
\toprule\noalign{}
Dimension & 说明 \\
\midrule\noalign{}
\endhead
\bottomrule\noalign{}
\endlastfoot
\textbf{核心Contribution} & State条件Data质量Framework + MLIP Verification \\
\textbf{Goal期刊} & npj Computational Materials / Nature
Communications \\
\textbf{受众} & MLIP 社区 + Data估值社区 \\
\textbf{篇幅} & 正文 1/3 理论（精简版）+ 2/3 Experiment \\
\end{longtable}

\paragraph{章节结构}\label{ux7ae0ux8282ux7ed3ux6784-1}

\begin{verbatim}
1. Introduction
   - MLIP Data贵，但Data价值不是固定的
   - 同一个势Function在不同构型区域性能差异巨大（引用 EGP P1 的 9.7× per-batch 差异）
   - 引出核心Problem：如何State条件地评估Data质量和专家可靠性？

2. Theory: State-Conditioned Expertise（精简到 ~3 页正文）
   2.1 Core Definitions
       R_m(s), SCX_m(s), Data四分类
   2.2 Three Propositions（精简版）
       Prop 1 (重构): 全局排序的遗憾下界
       Prop 2 (重构): 噪声存在时 error-driven sampling 的次优性
       Prop 3: State条件权重的优势
   2.3 Non-Identifiability Theorem
       无额外Assumption时，噪声与可学习困难不能仅从 (x,y) 识别
       → 多专家信息是必要的
   2.4 State Discovery: Two-Layer Architecture
       L1: 域知识Encoder → L2: ErrorDrivenEncoder（ErrorCorrelation子空间细化）

3. Methods: SCX Pipeline for MLIP
   3.1 ACE Descriptor State Encoder
   3.2 Two-Layer State Discovery
   3.3 State-Level Data Classification
   3.4 Action Policy

4. Experiments
   4.1 AlN v3: Two-Layer State Discovery（核心）
       - 一层 vs 两层对比 (F1: 0.253 → 0.585, Top-3: 37.8% → 94.6%)
       - 噪声Distribution（100% 在 thermal + MLMD）
       - Phonon 冗余（50-80% 可压缩）
   4.2 AlN v3: Retraining Validation ★P0 必做★
       - SCX 去噪 → 重训 → 实测力 RMSE vs baseline
       - vs random drop / loss-based drop
       - Training fmax vs test error: r=0.966
   4.3 Cross-Material Validation
       - Si/Cu/MgO 公开Data集上的 SCX Analysis
   4.4 Synthetic Experiments
       - 理论边界Verification（不可识别性Theorem的数值演示）

5. Discussion
   - SCX 是DataDiagnosis工具，与ModelArchitecture改进互补
   - 拓展潜力：工程仿真、医学Data
   - 局限：冷启动需要初始标签; State划分的粒度敏感
\end{verbatim}

\paragraph{不放入论文的内容}\label{ux4e0dux653eux5165ux8bbaux6587ux7684ux5185ux5bb9}

\begin{longtable}[]{@{}
  >{\raggedright\arraybackslash}p{(\linewidth - 2\tabcolsep) * \real{0.5714}}
  >{\raggedright\arraybackslash}p{(\linewidth - 2\tabcolsep) * \real{0.4286}}@{}}
\toprule\noalign{}
\begin{minipage}[b]{\linewidth}\raggedright
删除项
\end{minipage} & \begin{minipage}[b]{\linewidth}\raggedright
理由
\end{minipage} \\
\midrule\noalign{}
\endhead
\bottomrule\noalign{}
\endlastfoot
V(s) 乘法分解 & 启发式Definition，用四分类规则替代即可 \\
MedMNIST / CIFAR Experiment & 当前太弱，伤害可信度 \\
``六Dimension覆盖''叙事 & 改为''校准→路由→审计''三层 \\
Arrow 不可能Theorem类比 & 技术错误 + 过度包装 \\
LLM / Music / Semi / Vision / Sim Domain提及 & 收窄到 MLIP \\
``Data防Poisoning''概念 & 改为''Data质量审计'' \\
``general framework'' 声称 & 改为 ``methodology for MLIP data quality
with cross-domain potential'' \\
\end{longtable}

\begin{center}\rule{0.5\linewidth}{0.5pt}\end{center}

\subsubsection{3.4 第三跳：SCX-Sim 或
SCX-Health（选一个深做）}\label{ux7b2cux4e09ux8df3scx-sim-ux6216-scx-healthux9009ux4e00ux4e2aux6df1ux505a}

\begin{longtable}[]{@{}
  >{\raggedright\arraybackslash}p{(\linewidth - 6\tabcolsep) * \real{0.1875}}
  >{\raggedright\arraybackslash}p{(\linewidth - 6\tabcolsep) * \real{0.2812}}
  >{\raggedright\arraybackslash}p{(\linewidth - 6\tabcolsep) * \real{0.2812}}
  >{\raggedright\arraybackslash}p{(\linewidth - 6\tabcolsep) * \real{0.2500}}@{}}
\toprule\noalign{}
\begin{minipage}[b]{\linewidth}\raggedright
选项
\end{minipage} & \begin{minipage}[b]{\linewidth}\raggedright
Goal期刊
\end{minipage} & \begin{minipage}[b]{\linewidth}\raggedright
前提条件
\end{minipage} & \begin{minipage}[b]{\linewidth}\raggedright
优先级
\end{minipage} \\
\midrule\noalign{}
\endhead
\bottomrule\noalign{}
\endlastfoot
\textbf{SCX-Sim} & Nature Computational Science & 需产业合作方Data &
如有合作方则优先 \\
\textbf{SCX-Health} & npj Digital Medicine & MedMNIST Experiment需 GPU
升级且Result有效 & 如无合作方则选此项 \\
\end{longtable}

\textbf{建议二选一，不做两个。} 深度做透一个Domain比广撒网有效 10 倍。

\begin{center}\rule{0.5\linewidth}{0.5pt}\end{center}

\subsection{四、修改后的论文序列}\label{ux56dbux4feeux6539ux540eux7684ux8bbaux6587ux5e8fux5217}

\begin{longtable}[]{@{}
  >{\raggedright\arraybackslash}p{(\linewidth - 8\tabcolsep) * \real{0.1500}}
  >{\raggedright\arraybackslash}p{(\linewidth - 8\tabcolsep) * \real{0.1500}}
  >{\raggedright\arraybackslash}p{(\linewidth - 8\tabcolsep) * \real{0.2250}}
  >{\raggedright\arraybackslash}p{(\linewidth - 8\tabcolsep) * \real{0.2500}}
  >{\raggedright\arraybackslash}p{(\linewidth - 8\tabcolsep) * \real{0.2250}}@{}}
\toprule\noalign{}
\begin{minipage}[b]{\linewidth}\raggedright
顺序
\end{minipage} & \begin{minipage}[b]{\linewidth}\raggedright
论文
\end{minipage} & \begin{minipage}[b]{\linewidth}\raggedright
核心Problem
\end{minipage} & \begin{minipage}[b]{\linewidth}\raggedright
Goal期刊
\end{minipage} & \begin{minipage}[b]{\linewidth}\raggedright
DataState
\end{minipage} \\
\midrule\noalign{}
\endhead
\bottomrule\noalign{}
\endlastfoot
\textbf{1} & \textbf{EGP: Gauge-fixed shared+correction ACE} & ACE
势Function gauge fixing 与合并 & npj CompMat / PRM & ✅ AlN v3
已Training，GaN/AlN v4 超算中 \\
\textbf{2} & \textbf{SCX-MLIP: State-Conditioned Data Quality for MLIP}
& State条件Data质量Framework + MLIP Verification（理论 1/3 + Experiment 2/3） & npj CompMat
/ NatComm & ⏳ 需重训Verification + 公开Data集 \\
\textbf{3} & \textbf{SCX-Sim 或 SCX-Health} & 跨DomainExtension & NatCompSci /
npj DigMed & ⏳ 需合作方Data或 GPU Experiment \\
\end{longtable}

变化：原来的 SCX-Theory（单独投稿）和 SCX-MLIP 合并为一篇。SCX-Sim 和
SCX-Health 二选一。

\begin{center}\rule{0.5\linewidth}{0.5pt}\end{center}

\subsection{五、提交前必须完成的Experiment清单}\label{ux4e94ux63d0ux4ea4ux524dux5fc5ux987bux5b8cux6210ux7684ux5b9eux9a8cux6e05ux5355}

\begin{longtable}[]{@{}
  >{\raggedright\arraybackslash}p{(\linewidth - 6\tabcolsep) * \real{0.2222}}
  >{\raggedright\arraybackslash}p{(\linewidth - 6\tabcolsep) * \real{0.1667}}
  >{\raggedright\arraybackslash}p{(\linewidth - 6\tabcolsep) * \real{0.2500}}
  >{\raggedright\arraybackslash}p{(\linewidth - 6\tabcolsep) * \real{0.3611}}@{}}
\toprule\noalign{}
\begin{minipage}[b]{\linewidth}\raggedright
优先级
\end{minipage} & \begin{minipage}[b]{\linewidth}\raggedright
Experiment
\end{minipage} & \begin{minipage}[b]{\linewidth}\raggedright
计算成本
\end{minipage} & \begin{minipage}[b]{\linewidth}\raggedright
没有它的后果
\end{minipage} \\
\midrule\noalign{}
\endhead
\bottomrule\noalign{}
\endlastfoot
\textbf{P0} & \textbf{AlN v3 去噪后重训}：移除 fmax\textgreater5 帧或
SCX Top-2 噪声State → 重训 ACE → 测力 RMSE & 4-8h CPU &
审稿人直接拒稿------``预估不是Experiment'' \\
\textbf{P1} & \textbf{公开 MLIP Data集交叉Verification}：下载 Si/Cu/MgO
公开TrainingData → 运行 SCX 两层Analysis → Report噪声/冗余发现 & 1-2d Analysis（零
DFT） & ``只在 AlN 上有效''的批评无法回应 \\
\textbf{P1} & \textbf{合成Experiment}：生成已知噪声结构的合成Data → Verification SCX
恢复State结构的能力 → Verification不可识别性Theorem的边界 & 0.5d &
理论论文缺少对照Experiment \\
\textbf{P2} & \textbf{与 random drop / loss-based drop
的对比}：去噪Experiment中增加对照组 → 展示 SCX 指导的去噪优于 naive drop &
包含在 P0 中 & ``去掉高Error帧谁都会''的批评无法回应 \\
\textbf{P3} & \textbf{CIFAR-10 噪声检测 GPU 升级}：ResNet-18 + 50 epoch
+ 5 seeds & 1-2d GPU & 如果要做跨Domain claim 则需要；如果收窄到 MLIP
则不需要 \\
\end{longtable}

\begin{quote}
\textbf{P0 不可跳过。} 这个Experiment决定了论文能从 ``interesting
observation'' 升级到 ``verified contribution''。
\end{quote}

\begin{center}\rule{0.5\linewidth}{0.5pt}\end{center}

\subsection{六、理论加固清单}\label{ux516dux7406ux8bbaux52a0ux56faux6e05ux5355}

\begin{longtable}[]{@{}
  >{\raggedright\arraybackslash}p{(\linewidth - 4\tabcolsep) * \real{0.4000}}
  >{\raggedright\arraybackslash}p{(\linewidth - 4\tabcolsep) * \real{0.3000}}
  >{\raggedright\arraybackslash}p{(\linewidth - 4\tabcolsep) * \real{0.3000}}@{}}
\toprule\noalign{}
\begin{minipage}[b]{\linewidth}\raggedright
优先级
\end{minipage} & \begin{minipage}[b]{\linewidth}\raggedright
修改
\end{minipage} & \begin{minipage}[b]{\linewidth}\raggedright
说明
\end{minipage} \\
\midrule\noalign{}
\endhead
\bottomrule\noalign{}
\endlastfoot
\textbf{P0} & 删除 Arrow 不可能Theorem类比 & 技术错误，过度包装 \\
\textbf{P0} & 重构 Prop 1 &
从''不存在全局最优''改为''全局排序的遗憾下界为 Ω(p·δ\_min)'' \\
\textbf{P1} & 添加不可识别性Theorem &
显式Proof无额外Assumption时噪声与可学习困难不可区分 → 多专家信息的必要性 \\
\textbf{P1} & 添加冷启动协议 & Two-stage: 第一阶段无标签代理估计 →
第二阶段 acquisition \\
\textbf{P2} & V(s) 替换为四分类规则 & 避免启发式乘法分解争议 \\
\textbf{P2} & State发现抽象化 & 理论论文中将State映射视为
oracle，ErrorDrivenEncoder 作为Appendix实例化 \\
\end{longtable}

\begin{center}\rule{0.5\linewidth}{0.5pt}\end{center}

\subsection{七、CompetitionPositioning改进}\label{ux4e03ux7adeux4e89ux5b9aux4f4dux6539ux8fdb}

\begin{longtable}[]{@{}
  >{\raggedright\arraybackslash}p{(\linewidth - 4\tabcolsep) * \real{0.2308}}
  >{\raggedright\arraybackslash}p{(\linewidth - 4\tabcolsep) * \real{0.3462}}
  >{\raggedright\arraybackslash}p{(\linewidth - 4\tabcolsep) * \real{0.4231}}@{}}
\toprule\noalign{}
\begin{minipage}[b]{\linewidth}\raggedright
对手
\end{minipage} & \begin{minipage}[b]{\linewidth}\raggedright
当前Positioning
\end{minipage} & \begin{minipage}[b]{\linewidth}\raggedright
改进后Positioning
\end{minipage} \\
\midrule\noalign{}
\endhead
\bottomrule\noalign{}
\endlastfoot
\textbf{Data Shapley} & 替代关系（SCX V(s) vs Data Shapley φ\_i） &
\textbf{互补关系}：SCX 粗估State → Data Shapley 在选定State内细估 \\
\textbf{MoE (Shazeer/DeepMD)} & 两阶段 vs 端到端 &
\textbf{场景差异}：异构专家 + 增量Extension + 可审计 → SCX；同构端到端 → 传统
MoE \\
\textbf{LLM Data治理} & 不做主战场 & \textbf{Introduction一句话 +
立刻转移}：先提 LLM 场景作为动机，然后转移到科学仿真（有 gold label） \\
\textbf{条件风险} & 新概念 & \textbf{视角差异}：离散化State ≠
连续条件风险；离散化是审计的前提 \\
\end{longtable}

\begin{center}\rule{0.5\linewidth}{0.5pt}\end{center}

\subsection{八、论文间叙事关系}\label{ux516bux8bbaux6587ux95f4ux53d9ux4e8bux5173ux7cfb}

\begin{verbatim}
EGP Paper 1 (concrete, 材料Method)
    │
    │ "我们在这个工作中发现 expert merging 需要 gauge fixing + state awareness"
    │   → 为 SCX-MLIP 提供经验动机
    │
    ▼
SCX-MLIP (theory + application, 合并版)
    │
    │ "EGP P1 的经验观察 → 形式化为State条件专家性理论 → 回到 MLIP Verification → 跨材料Verification"
    │  理论 1/3 篇幅 + Experiment 2/3 篇幅
    │
    ▼
SCX-Sim 或 SCX-Health (cross-domain)
    │ "SCX Framework在另一个完全不同的Domain的适用性"
\end{verbatim}

\subsubsection{引用关系}\label{ux5f15ux7528ux5173ux7cfb}

\begin{verbatim}
EGP P1 ──(不引用 SCX)──▶ Independent发表
SCX-MLIP ──(引用 EGP P1 作为 motivation)──▶ 在 EGP 发表后投稿
SCX-Sim/Health ──(引用 SCX-MLIP 作为Method论基础)──▶ 在 SCX-MLIP 发表后投稿
\end{verbatim}

\begin{center}\rule{0.5\linewidth}{0.5pt}\end{center}

\subsection{九、一句话总结}\label{ux4e5dux4e00ux53e5ux8bddux603bux7ed3}

\textbf{``一击即中''的Core不是理论有多漂亮，而是审稿人读完第三页的时候已经被
AlN 的 9.7 倍 per-batch
差异和去噪重训的实测改善说服了------此时你再告诉他背后的理论Framework，他只会觉得''这个理论有道理，因为它解释了我刚才看到的Experiment现象''。}

\begin{center}\rule{0.5\linewidth}{0.5pt}\end{center}

\emph{本Report由 5 个Independent AI Review员并行Review后Synthesis生成。各Review员IndependentReport见
CodexKnowledge/agent\_outputs/ 目录。}

\end{document}
